\section{Thermal Fin Benchmark}
\label{sec:results_thermal_fin}

The Thermal Fin benchmark is the final and most demanding case in this study.
It originates from a 2D steady-state parametric problem in the RBniCS
library~\cite{rozza2008}.
We adapted it to a full 3D transient water-quenching scenario, making it the
most geometrically complex domain in the evaluation.
The geometry consists of a rectangular base plate with four cooling fins.
We evaluate the transient temperature field against the FEM reference solution
produced with a Crank--Nicolson solver.
The acceptance threshold is $15\,^\circ\mathrm{C}$ mean-window MAE,
corresponding to roughly 3\% of the full operating temperature range.

Two error values are reported throughout this section. The
\emph{mean-window MAE} is the main engineering metric. It averages the
prediction error over the complete 30-second rollout. The \emph{final window
MAE} is also included when discussing the plotted temperature fields at
$t=30\,\mathrm{s}$. Readers should not treat it as a primary indicator. The
field becomes nearly uniform by the end of quenching, so a low final-window
value does not mean the network was accurate throughout the simulation.

Table~\ref{tab:tf_overview} collects the headline numbers from each
experimental mode before we go into the details.

\begin{table}[H]
\centering
\caption{Thermal Fin Benchmark: headline results across all experimental modes.
$^*$Bayesian/TPE uses Fourier feature embedding; NSGA-II/III use plain MLP.}
\label{tab:tf_overview}
\small
\begin{tabular}{llccc}
\toprule
Mode & Architecture & Acceptable $k$ & Best MAE ($^\circ\mathrm{C}$) & FEM saved \\
\midrule
FEM-anchored, fixed $k$ & Bayesian/TPE$^*$ & $1$--$5$ &  5.10 & up to 80\% \\
FEM-anchored, fixed $k$ & NSGA-II      & $1$--$4$ &  8.98 & up to 75\% \\
FEM-anchored, fixed $k$ & NSGA-III     & $1$--$4$ &  8.97 & up to 75\% \\
\midrule
FEM-anchored, adaptive  & Bayesian/TPE & —        & 11.26 & 45\% \\
FEM-anchored, adaptive  & NSGA-II      & —        & \textbf{10.54} & 45\% \\
FEM-anchored, adaptive  & NSGA-III     & —        & 13.21 & 45\% \\
\midrule
PINN-only, best $k$ & Bayesian/TPE & — (all fail) & 217.10 & 95\% \\
PINN-only, best $k$ & NSGA-II      & — (all fail) & 113.46 & 95\% \\
PINN-only, best $k$ & NSGA-III     & — (all fail) &  59.20 & 95\% \\
\bottomrule
\end{tabular}
\end{table}

Table~\ref{tab:tf_overview} serves as a roadmap for this section. The
\emph{Acceptable $k$} column identifies the highest skip factor within the
$15\,^\circ\mathrm{C}$ engineering acceptance band. Bayesian/TPE (Fourier)
reaches $k=5$ (80\% FEM saving). NSGA-II and NSGA-III top out at $k=4$
(75\%). Adaptive-$k$ control achieves 45\% saving for all three architectures
but does not exceed fixed-$k$ accuracy. PINN-only mode fails completely,
demonstrating that periodic FEM correction is not optional on this geometry.

% -----------------------------------------------------------------------------
\subsection{FEM-Anchored Fixed-$k$ Sweep}
\label{subsec:tf_fem_fixed}

In FEM-anchored mode the PINN predicts over a window of $k \times 1.5\,\mathrm{s}$
and is corrected by the next FEM snapshot.  Larger $k$ means fewer FEM calls
but a longer prediction interval.  All three architectures were trained at
every $k$ from 1 to 5 (15 configurations total), using 2\,000 Adam epochs
followed by 100 L-BFGS steps with 12\,000 interior and 2\,400 boundary
collocation points per window.

Full per-seed MAE and $L_2$ tables are in Appendix~\ref{app:full_results}
(Tables~\ref{tab:tf_consolidated} and~\ref{tab:tf_l2}).

Bayesian/TPE (Fourier) achieves $L_2{=}0.049$ at $k=1$, rising to $0.081$
at $k=5$, which is still well below the $L_2{=}0.10$ practical threshold.
NSGA-II and NSGA-III follow a similar trend, reaching $0.083$--$0.084$ at
$k=5$.
Seed variance is very low ($\sigma \leq 0.007$), confirming reproducibility.

Three key findings stand out. First, Bayesian/TPE with Fourier embedding remains within
the $15\,^\circ\mathrm{C}$ acceptance band at every skip factor including
$k=5$ ($12.55\pm0.30\,^\circ\mathrm{C}$, 80\% FEM saving). Four FEM calls
suffice instead of twenty. Seed variance is very low
($\sigma \leq 0.29\,^\circ\mathrm{C}$), confirming that this result is
reproducible. Second, NSGA-II and NSGA-III remain acceptable through $k=4$
($14.14$ and $14.40\,^\circ\mathrm{C}$ respectively, 75\% saving) but exceed
the threshold at $k=5$ by less than $1\,^\circ\mathrm{C}$. Third, Bayesian/TPE
consistently achieves the lowest mean at every $k$. It benefits from both its
larger capacity ($5{\times}151$ vs.\ $3{\times}153$ and $3{\times}75$) and the
Fourier embedding that reduces per-window error by $1.8$--$4.4\,^\circ\mathrm{C}$
relative to plain MLP (see Appendix~\ref{app:ablation} for the ablation).

Prediction quality degrades gradually as $k$ increases and errors concentrate at the fin-base junctions.
Temperature field heat maps (2D cross-sections and 3D volumetric views) for all architectures and skip factors are in Appendix~\ref{app:tf_heatmaps} (Figures~\ref{fig:app_tf_fixed_2d} and~\ref{fig:app_tf_fixed_3d}).

% -----------------------------------------------------------------------------
\subsection{FEM Anchored Adaptive \texorpdfstring{$k$}{k}}
\label{subsec:tf_fem_adaptive}

The adaptive controller monitors the per-window error and adjusts $k$ on the
fly. It starts conservatively at $k=1$ and promotes to a longer window only
once the error has stayed low for two consecutive windows. Phase-based caps
prevented promotion to $k>2$ before $t=10\,\mathrm{s}$ and to $k>4$ before
$t=20\,\mathrm{s}$, ensuring the controller kept short windows during the
steepest part of the transient. All three architectures were trained with the
same adaptive schedule, using the same per-window training budget as the
fixed-$k$ runs. Table~\ref{tab:tf_fem_adaptive} shows the outcomes.

\begin{table}[H]
\centering
\caption{Thermal Fin adaptive-$k$ controller: FEM savings, accuracy, and window schedule.}
\label{tab:tf_fem_adaptive}
\resizebox{\textwidth}{!}{%
\begin{tabular}{lcccccl}
\toprule
Architecture & FEM calls & FEM saved & Mean MAE ($^\circ$C) & $L_2$ rel. & Final MAE ($^\circ$C) & $k$ schedule \\
\midrule
Bayesian/TPE & 11 & 45\% & 11.26 & 0.068 & 6.35
  & 1-1-1-1-1-2-2-3-3-4-1 \\
NSGA-II      & 11 & 45\% & \textbf{10.54} & \textbf{0.061} & \textbf{4.50}
  & 1-1-1-1-1-2-2-3-3-4-1 \\
NSGA-III     & 11 & 45\% & 13.21 & 0.068 & 6.06
  & 1-1-1-1-1-1-2-2-3-3-4 \\
\bottomrule
\end{tabular}}
\end{table}

Table~\ref{tab:tf_fem_adaptive} shows MAE and $L_2$ relative error alongside
the window schedule selected by the controller. The controller used 11 of 20
FEM calls (45\% saving), and the remaining 9 FEM calls were replaced by PINN
predictions over progressively longer windows. The adaptive schedule does not
outperform the best fixed-$k$ result. Bayesian/TPE with fixed $k=5$ achieves
$12.55\,^\circ\mathrm{C}$ (ten-seed mean) at 80\% FEM saving, while the adaptive run reaches
only $11.26\,^\circ\mathrm{C}$ at 45\% saving. The mechanism behind this
underperformance, a threshold--architecture mismatch that causes different
architectures to promote and demote at incorrect points, is analysed in
Section~\ref{sec:disc_thermalfin}. The $k$ schedule column lists the skip
factor for each consecutive window from $t=0$ to $t=30\,\mathrm{s}$. The
final entry ($k=1$) covers the last short window near equilibrium. All three
architectures achieve mean-window MAE within the $15\,^\circ\mathrm{C}$
threshold and $L_2{<}0.07$ (excellent field reconstruction), confirming that
the adaptive controller maintains quality comparable to fixed $k=2$--$3$ runs.

The controller behaved physically as expected: it held at $k=1$ during the
steep early gradient phase, then promoted to $k=2$, $k=3$, and $k=4$ as the
field relaxed, automatically tracking prediction difficulty.
The window timing for each architecture is shown in the $k$ schedule column of
Table~\ref{tab:tf_fem_adaptive}.

Temperature field heat maps for the adaptive-$k$ runs are in Appendix~\ref{app:tf_heatmaps} (Figure~\ref{fig:app_tf_adap_pinn_2d}).

Figure~\ref{fig:tf_l2_curve} shows how the per-window $L_2$ relative error
evolves over the 30-second quench. Error is highest in the first windows
because fin-tip gradients are steepest early in the quench. The error
decreases as the field approaches uniformity, consistent with the early-window
difficulty seen in all benchmark groups.

\begin{figure}[H]
\centering
\includegraphics[width=\textwidth]{l2_all_k_thermal_fin.png}
\caption{Per-window $L_2$ relative error over the 30-second Thermal Fin quench.
Line: ten-seed mean; band: $\pm$one standard deviation. Colour encodes $k$.}
\label{fig:tf_l2_curve}
\end{figure}


% -----------------------------------------------------------------------------
\subsection{PINN Only Rollouts}
\label{subsec:tf_pinn_only}

To understand how much of the accuracy comes from the FEM corrections, the
network was also run without any anchoring after $t=0$. In this mode the
network receives the FEM initial condition once and then propagates entirely
from its own output, with no corrections at any subsequent window. This is the
most extreme test of the surrogate. Each window's prediction depends only on
the previous PINN output, so any error in one window feeds directly into the
next as its initial condition. All 15 architecture--$k$ combinations were run
under this protocol, with the same training setup as the FEM-anchored runs.
Table~\ref{tab:tf_pinn_only_fixed} shows the full set of results.

\begin{table}[H]
\centering
\caption{Thermal Fin PINN-Only Fixed-$k$ Rollout: Catastrophic Error Without FEM Correction.
Mean-window MAE and final-window MAE ($^\circ\mathrm{C}$) across all 15 architecture--skip-factor
combinations run without any FEM anchor after $t=0$.  Bold = lowest mean MAE per architecture.
$L_2$ relative error is not reported because all values exceed the practical measurement range ($L_2 \gg 1$).}
\label{tab:tf_pinn_only_fixed}
\small
\setlength{\tabcolsep}{4pt}
\begin{tabular}{lcccccc}
\toprule
Architecture & $k$ & Windows & FEM saved & Mean MAE ($^\circ$C) & Final MAE ($^\circ$C) \\
\midrule
Bayesian/TPE & 1 & 20 & 95\% & 217.10 & 327.97 \\
Bayesian/TPE & 2 & 10 & 95\% & 242.96 & 348.91 \\
Bayesian/TPE & 3 &  7 & 95\% & 249.30 & 345.20 \\
Bayesian/TPE & 4 &  5 & 95\% & 221.42 & 292.17 \\
Bayesian/TPE & 5 &  4 & 95\% & 264.88 & 350.28 \\
\midrule
NSGA-II      & 1 & 20 & 95\% & 180.46 & 241.67 \\
NSGA-II      & 2 & 10 & 95\% & 119.62 &  57.13 \\
NSGA-II      & 3 &  7 & 95\% & \textbf{113.46} &  50.85 \\
NSGA-II      & 4 &  5 & 95\% & 134.24 &  83.25 \\
NSGA-II      & 5 &  4 & 95\% & 160.15 & 108.26 \\
\midrule
NSGA-III     & 1 & 20 & 95\% &  83.78 & 115.15 \\
NSGA-III     & 2 & 10 & 95\% & 160.89 & 206.46 \\
NSGA-III     & 3 &  7 & 95\% &  62.90 &  51.16 \\
NSGA-III     & 4 &  5 & 95\% &  95.91 &  66.54 \\
NSGA-III     & 5 &  4 & 95\% & \textbf{59.20} &  44.47 \\
\bottomrule
\end{tabular}
\end{table}

Table~\ref{tab:tf_pinn_only_fixed} reports mean-window and final-window MAE
for every (architecture, $k$) combination without FEM anchoring.
\emph{Mean-window MAE} averages the absolute error over all prediction windows
and grid points throughout the 30-second rollout. This is the honest measure
of overall accuracy. \emph{Final-window MAE} measures only the last window at
$t=30\,\mathrm{s}$. Because the temperature field approaches uniformity at the
end of quenching, final-window values appear deceptively low even when
intermediate predictions were catastrophically wrong. Mean-window MAE should
always be used as the primary metric.

Bayesian/TPE is completely unstable without FEM correction. All five $k$
values produce mean-window MAE above $200\,^\circ\mathrm{C}$, showing that
the large ReLU network accumulates drift regardless of window length. NSGA-III
performs best (minimum $59.20\,^\circ\mathrm{C}$ at $k=5$), yet this is still
nearly four times the $15\,^\circ\mathrm{C}$ threshold. $L_2$ relative error
is not reported for PINN-only runs because all values exceed the practical
range ($L_2 \gg 1$), making the metric uninformative.

Unlike the FEM-anchored runs where errors stay localised near fin junctions,
the PINN-only errors spread over the entire fin geometry.
Temperature field heat maps for the PINN-only runs are in Appendix~\ref{app:tf_heatmaps} (Figure~\ref{fig:app_tf_adap_pinn_2d}).

The adaptive controller was also tested in PINN-only mode.
Table~\ref{tab:tf_pinn_only_adaptive} shows the outcomes.

\begin{table}[H]
\centering
\caption{Thermal Fin PINN-Only Adaptive-$k$ Rollout: Adaptive Controller Failure Without FEM Correction.
The controller stalls at $k=1$ for all architectures because the first-window error immediately
exceeds the demotion threshold, preventing any window promotion.  Mean-window MAE exceeds $100\,^\circ\mathrm{C}$
in all cases, demonstrating that periodic FEM correction is a necessary condition for stable prediction.}
\label{tab:tf_pinn_only_adaptive}
\resizebox{\textwidth}{!}{%
\begin{tabular}{lccccl}
\toprule
Architecture & FEM calls & FEM saved & Mean MAE ($^\circ$C) & Final MAE ($^\circ$C) & $k$ schedule \\
\midrule
Bayesian/TPE & 1 & 95\% & 200.40 & 199.70 & $k{=}1$ for all 20 windows \\
NSGA-II      & 1 & 95\% & 173.62 & 185.99 & $k{=}1$ for all 20 windows \\
NSGA-III     & 1 & 95\% & 185.00 & 197.14 & $k{=}1$ for all 20 windows \\
\bottomrule
\end{tabular}}
\end{table}

Table~\ref{tab:tf_pinn_only_adaptive} shows the extreme consequence of
removing FEM anchoring entirely. The adaptive controller evaluates whether to
increase $k$ based on per-window error. Without FEM correction, the
first-window error exceeds $100\,^\circ\mathrm{C}$ immediately, above the
demotion threshold, so the controller stalls at $k=1$ for all 20 windows.
Even maximum temporal resolution (no step-skipping) cannot save the rollout.
The error comes from the absence of periodic FEM correction, not from window
length. This table is the clearest demonstration in the thesis that FEM
anchoring is \emph{necessary}, not merely helpful.

Figure~\ref{fig:tf_seed_bars} shows the FEM-anchored fixed-$k$ mean-window
MAE for all ten seeds across each skip factor. Each panel corresponds to one
architecture. The solid line is the ten-seed mean and the shaded band covers
one standard deviation.

\begin{figure}[H]
\centering
\includegraphics[width=\textwidth]{seed_bar_thermal_fin.png}
\caption{Seed-to-seed variability in mean-window MAE for the Thermal Fin.
Line: ten-seed mean; band: $\pm$one standard deviation.}
\label{fig:tf_seed_bars}
\end{figure}

Three patterns stand out. First, Bayesian/TPE (Fourier) shows the narrowest
bands at every skip factor, confirming that the Fourier embedding stabilises
training against random initialisation as well as improving accuracy. Second,
the band width grows with $k$ for all architectures, consistent with the
observation that longer prediction windows amplify initialisation-dependent
variability. Third, no individual seed violates the $15\,^\circ\mathrm{C}$
threshold for Bayesian/TPE through $k=5$, whereas NSGA-II and NSGA-III have
individual seeds above the threshold only at $k=5$ (see
Table~\ref{tab:tf_consolidated} for the full per-seed values).

The low seed variance across the thermal fin domain, which is smaller than for
most 3D canonical domains, reflects the effect of the Fourier embedding for
Bayesian/TPE and the relatively moderate geometry complexity for NSGA-II/III.
The fin junction gradients are steep but localised, and the MSWP strategy
limits error accumulation through periodic FEM correction.
